\scanpage{21}\mspage{20}
pour tout $A\in F$, $B\in G$, $A\cap F$ est \deleted{réunion de}\inserted{une} partie $F$-fermée de $|F|$, i.e. réunion de strates de $F$.
% slanted
\margin{\uline{NB} Le ``de plus'' implique déjà la première condition, si on ajoute que $|\deleted{\unsure{A}}F|\subset|\deleted{\unsure{B}}G|$, et cela équivaut à $|F|\unsure{-|\longrightarrow}|G|$ \uline{continue} (ou croissante).}\margin{Cf condition équivalente $\mathrm{b}'$) à la pge suivante}

b) Pour tout $A\in F$, $\exists\,B\in G$ telle que $A^\circ\subset B^\circ$ (strate évidemt unique, \inserted{\unsure{noté|notée|soit} $\varphi(A)$}), l'application $\varphi:\deleted{\unsure{F}}F\unsure{-|\longrightarrow}G$ ainsi obtenue est croissante (
\deleted{$A'\subset A$, et} \deleted{$A'^\circ\in\unsure{+|F}$} \deleted{$A'\subset A$} $A,\deleted{\unsure{B}}A'\in F$, $\deleted{\unsure{A'}}B,B'\in G$, \deleted{$A\unsure{\in|=|\subset} A'$, $A^\circ\subset B^\circ$, $\unsure{B\subset B'}$} avec $A^\circ\subset B^\circ$, $A'^\circ\subset B'^\circ$, $A'\subset A$, $\Longrightarrow B'\subset B$.

a) $\Longrightarrow$ b) Supposons a) prouvons b). Soit $A\in F$, \deleted{\unsure{par}} a) $\exists\ B\inserted{{}=\varphi(A)}\in G$ \deleted{tel} qui soit plus petit ds $G^A$. Prouvons $\deleted{\unsure{B}}A^\circ\subset B^\circ$. Soit
$x\in A^\circ$, il faut prouver que $B'\in G$, $B'\deleted{\unsure{\in}}\unsure{\subset|\not\subset} B$ et $x\in B'\Longrightarrow B'=B$. \deleted{Soit} Considérons $B'\cap A$, c'est une partie de $A$ réunion \unsure{de'|de} strates de $F$,
donc $\exists\ A'\in\deleted{\unsure{\Phi}}F$, telle que $x\in A'\subset B'\cap A$. Comme $x\in A^\circ$ et $A'\in\deleted{\unsure{\Phi}}F_A$, on a $A'=A$ et comme $A'=A\subset B'\cap A$, on a $A\subset B'$. D'après le caractère minimal de
$B^{\unsure{s|'}}$, on a donc $B'=B$ qed. / D'autre part, on a déjà noté que l'application $\varphi$ est croissante. \deleted{\unsure{Donc}}Cela établit b).

b) $\Longrightarrow$ a). Supposons b, prouvons a). Soit $A\in F$, soit $B\in G$ telle que $A^\circ\subset B^\circ$, prouvons que $B$ est plus petit \inserted{élémt} dans $G^A$.
% slanted
\margin{\uline{Lemme} Moyennant b) on a \unsure{(b)} : Si $A\in F$, $B\in G$, $A^\circ\subset B$ alors $A\subset B$}
Prouvons déjà que $B\in G^A$, i.e. $A\subset B$ \deleted{Soit $B'\in G$ tel que $A\subset B'$}
\deleted{Soit} On a $A=\deleted{\bigcup_{\unsure{A'}}\unsure{A^\circ}\unsure{\in|\subset}\unsure{\sim}}\bigcup_{\substack{A'\in F\\ A'\unsure{\subset|\in} A}} A'^\circ$, \deleted{Comme $A^\circ\subset B^\circ\subset B$,} il suffit de prouver que les $A'^\circ$ sont $\subset B$
